\begin{abstract}
视觉语言模型把一幅图像编码成很长的 patch 序列。以 LLaVA-1.5 的常用设置为准，576 个视觉 token 进入 32 层 Decoder，既增加预填充中的注意力计算，也拉长后续解码要读取的缓存。把序列压到很小以后，选择不再只是丢掉明显背景：补进与当前问题更相关的位置，可能同时拆掉一批原本由它代表的视觉位置。本文提出 BELA。它先用视觉塔两个深度的特征做一个与问题无关的固定容量初始化，再允许问题信号对这个集合做有限次一进一出的交换；一次交换能否被接受，取决于新位置带来的问题增益，以及全体原始位置上最近代表覆盖下降了多少。被选中的仍是原模型送入投影器的特征。进入 Decoder 之后，这些 token 先参与浅层计算，再按预定层数逐段减少。在 LLaVA-1.5-7B 上，层平均保留 32 个视觉 token，对应的视觉 token-layer 预算是完整视觉序列的 $1/18$。与同一日程、同一种后层规则的原生集合相比，问题条件交换把六项任务相对完整模型的平均比值从 93.56\% 提高到 95.34\%。GQA 与 POPE F1 上升，TextVQA-OCR 下降。该比值不是答题正确率，预算比例也不是实测加速。
\end{abstract}
